\documentclass[12pt]{article}
\usepackage[utf8]{inputenc}
\usepackage[T1]{fontenc}
\usepackage[margin=1in]{geometry}
\usepackage{setspace}
\usepackage{xcolor}
\usepackage[colorlinks=true, urlcolor=blue!60!black, linkcolor=blue!60!black]{hyperref}
\setlength{\parindent}{0em}
\setlength{\parskip}{0.6em}
\renewcommand{\baselinestretch}{1.0}
\begin{document}
\begin{center}
{\LARGE \textit{Virginia State Board of Pharmacy v. Virginia Citizens Consumer Council}, 425 U.S. 748 (1976)\par}
\vspace{0.5em}
{\large Case Brief --- Two-Page Edition\par}
\end{center}

\section{Memory Jogger}
Commercial speech is protected: Virginia's ban on advertising prescription-drug prices fell, and consumers may assert a right to receive the information (\href{https://supreme.justia.com/cases/federal/us/425/748/}{425 U.S.\ 748}).

\section{Detailed Case Facts}
Virginia made it ``unprofessional conduct'' for pharmacists to advertise prescription-drug prices, effectively banning all such advertising. Consumers and consumer groups sued, stipulating that prices for the same drugs varied by as much as 650\% within a single area and that pharmacists would advertise if allowed.

\section{Procedural History}
A three-judge district court held the ban unconstitutional; the Board appealed directly.

\section{Judicial Votes}
Affirmed, 7--1. Blackmun wrote for Burger, Brennan, Stewart, White, Marshall, and Powell; Burger and Stewart concurred; Rehnquist dissented; Stevens did not participate.

\section{Holding}
Consumers have a reciprocal right to receive truthful advertising from willing speakers. Purely commercial speech is protected, and a State may not suppress truthful information about lawful activity out of fear of its effects on those who receive it.

\section{Analysis of Opinions}
Blackmun rejected \textit{Valentine v.\ Chrestensen}'s exception: speech that does ``no more than propose a commercial transaction'' serves the consumer, especially ``the poor, the sick, and particularly the aged,'' and society's interest in informed economic decisions. Virginia's ``highly paternalistic approach'' kept the public ignorant instead of regulating pharmacists directly. The State may still bar false or misleading ads and impose time, place, and manner rules. Burger reserved the learned professions; Stewart stressed that false ads remain regulable. Rehnquist accused the Court of reviving \textit{Lochner} and predicted ads for lawyers, liquor, and drugs like Demerol.

\section{Examples: Future Applications}
\textbf{Same-Side:} the First Amendment protects lawyers who advertise flat fees for routine services (\textit{Bates}). \textbf{Opposite-Side:} in-person solicitation of accident victims may be banned (\textit{Ohralik}). \textbf{Fence-Sitter:} a ban on sportsbook signup-bonus advertising.

\section{Critique}
Critics say Blackmun's link to self-government constitutionalizes laissez-faire assumptions. Rehnquist's dissent anticipated direct-to-consumer drug advertising.

\section{Key Quotations}
\begin{quote}
``As to the particular consumer's interest in the free flow of commercial information, that interest may be as keen, if not keener by far, than his interest in the day's most urgent political debate. \dots\ So long as we preserve a predominantly free enterprise economy, the allocation of our resources in large measure will be made through numerous private economic decisions. It is a matter of public interest that those decisions, in the aggregate, be intelligent and well informed. To this end, the free flow of commercial information is indispensable.''
\end{quote}
\noindent --- Justice Blackmun, 425 U.S.\ at 763, 765.

\end{document}
